% ============================================================================
% 05_fusion_framework.tex — MSRF 独立方法章（新 §5）
% 新版提示词 P1：将 MSRF 提升为独立方法章节，形式化 + 算法框 + 架构图。
% 全部数字与结构取自 stage_p2_msrf.py 与结果产物（见 REVISION_LOG R28）。
% ============================================================================
\section{A Knowledge-Driven Multi-Source Fusion Framework (MSRF)}
\label{sec:msrf}

Single-cue detectors are not trustworthy on framed inputs: off-the-shelf
harmfulness classifiers fail catastrophically on Chinese inputs (FNR
$\geq 0.96$; Section~\ref{sec:setup}) and the strongest open guards score
below chance on the framing distribution (Section~\ref{sec:fusion}). The
preceding analysis also establishes that the effect is carried by a specific,
identifiable component --- narrative structure --- not by generic ``attackness''
or modality. \MSRF{} is a knowledge-driven response to both facts: it
represents the evidence about a response as \emph{four mutually
non-overlapping knowledge sources}, calibrates each source's score, and fuses
them under an explicit uncertainty treatment into a single auditable score.
Figure~\ref{fig:msrf-framework} summarizes the architecture.

\begin{figure}[t]
\centering
\begin{tikzpicture}[
  font=\footnotesize,
  bb/.style={rounded corners=3pt, draw=gray!45, dashed, fill=figBg},
  lbl/.style={font=\scriptsize\bfseries, text=gray!60},
  flw/.style={font=\scriptsize, text=gray!65},
  box/.style={rounded corners=2pt, draw=figGray, fill=white, inner sep=3pt},
  main/.style={-{Stealth[length=2.2mm]}, draw=gray!70, line width=1.1pt},
  aux/.style={-{Stealth[length=1.8mm]}, draw=gray!60, line width=0.8pt},
  bus/.style={draw=gray!70, line width=1.1pt},
]
% ---- swimlane bands (labels anchored top-left in reserved strip, inner margin >= 14pt) ----
\draw[bb] (0.08,0.12) rectangle (2.32,7.70);
\draw[bb] (2.42,0.12) rectangle (6.30,7.70);
\draw[bb] (6.40,0.12) rectangle (10.66,7.70);
\draw[bb] (10.76,0.12) rectangle (13.32,7.70);
\node[lbl, anchor=west] at (0.23,6.95) {INPUT};
\node[lbl, anchor=west] at (2.57,6.95) {KNOWLEDGE};
\node[lbl, anchor=west] at (6.62,6.95) {CALIB \& FUSION};
\node[lbl, anchor=west] at (10.91,6.95) {DECISION};

% ---- input: LALM response (title above box, >=6pt; text/audio dual channel) ----
\node[box, rounded corners=3pt, minimum width=2.2cm, minimum height=0.7cm] at (1.2,3.30) {};
\node[font=\scriptsize\bfseries] at (1.2,4.10) {LALM resp.};
\fill[intentC!10] (0.12,2.95) rectangle (1.20,3.55);
\fill[gray!10]    (1.20,2.95) rectangle (2.28,3.55);
\draw[gray!60]    (1.20,2.95) -- (1.20,3.55);
\node[font=\footnotesize] at (0.66,3.25) {text};
\node[font=\footnotesize] at (1.74,3.25) {audio};

% ---- knowledge sources (semantic colors + feature dims from Eqs.~1--4; d=3/4/5/1) ----
% Intent (blue, d=3)
\node[box, rounded corners=3pt, minimum width=3.35cm, minimum height=1.0cm,
      fill=intentC!12, draw=intentC] at (4.6,6.05) {};
\fill[intentC] (2.94,5.67) rectangle ++(0.14,0.76);
\node[right, font=\footnotesize\bfseries, text=intentC!85!black] at (3.16,6.30) {Intent};
\node[right, font=\scriptsize, text=black!85] at (3.16,5.78) {LoRA prob. (d=3)};
% Narrative (orange, d=4)
\node[box, rounded corners=3pt, minimum width=3.35cm, minimum height=1.0cm,
      fill=narrativeC!12, draw=narrativeC] at (4.6,4.80) {};
\fill[narrativeC] (2.94,4.42) rectangle ++(0.14,0.76);
\node[right, font=\footnotesize\bfseries, text=narrativeC!85!black] at (3.16,5.05) {Narrative};
\node[right, font=\scriptsize, text=black!85] at (3.16,4.53) {discourse feats (d=4)};
% Acoustic (green, d=5)
\node[box, rounded corners=3pt, minimum width=3.35cm, minimum height=1.0cm,
      fill=acousticC!12, draw=acousticC] at (4.6,3.55) {};
\fill[acousticC] (2.94,3.17) rectangle ++(0.14,0.76);
\node[right, font=\footnotesize\bfseries, text=acousticC!85!black] at (3.16,3.80) {Acoustic};
\node[right, font=\scriptsize, text=black!85] at (3.16,3.28) {prosody (d=5)};
% Uncertainty (purple, d=1)
\node[box, rounded corners=3pt, minimum width=3.35cm, minimum height=1.0cm,
      fill=uncertaintyC!12, draw=uncertaintyC] at (4.6,2.30) {};
\fill[uncertaintyC] (2.94,1.92) rectangle ++(0.14,0.76);
\node[right, font=\footnotesize\bfseries, text=uncertaintyC!85!black] at (3.16,2.55) {Uncertainty};
\node[right, font=\scriptsize, text=black!85] at (3.16,2.03) {disagreement (d=1)};

% ---- fusion chain ----
\node[box, rounded corners=3pt, minimum width=2.4cm, minimum height=1.0cm] at (9.35,6.05) {};
\node[font=\footnotesize\bfseries] at (9.35,6.30) {out-of-fold};
\node[font=\scriptsize, text=black!80] at (9.35,5.78) {group-aware cv};
\node[box, rounded corners=3pt, minimum width=2.4cm, minimum height=1.0cm] at (9.35,4.85) {};
\node[font=\footnotesize\bfseries] at (9.35,5.10) {isotonic};
\node[font=\scriptsize, text=black!80] at (9.35,4.58) {$\sigma_i$ on train folds};
% fusion MLP (title as first in-box line; 4 stacked layer boxes 4->16->8->1)
\node[box, rounded corners=3pt, minimum width=3.3cm, minimum height=2.7cm] at (8.9,2.55) {};
\node[font=\footnotesize\bfseries] at (8.9,3.55) {225-param MLP};
% layer 1: input (4) — left edge carries the four source colours
\draw[rounded corners=1pt, fill=white, draw=gray!70, line width=1pt]
  (8.05,2.86) rectangle (9.65,3.24);
\fill[intentC]      (8.18,3.145) rectangle (8.30,3.24);
\fill[narrativeC]   (8.18,3.05)  rectangle (8.30,3.145);
\fill[acousticC]    (8.18,2.955) rectangle (8.30,3.05);
\fill[uncertaintyC] (8.18,2.86)  rectangle (8.30,2.955);
\node[font=\scriptsize] at (8.85,3.05) {\textbf{4} input};
% layer 2: hidden 16 (ReLU)
\draw[rounded corners=1pt, fill=white, draw=gray!70, line width=1pt]
  (8.05,2.36) rectangle (9.65,2.74);
\node[font=\scriptsize] at (8.85,2.55) {\textbf{16} ReLU};
% layer 3: hidden 8 (ReLU)
\draw[rounded corners=1pt, fill=white, draw=gray!70, line width=1pt]
  (8.05,1.86) rectangle (9.65,2.24);
\node[font=\scriptsize] at (8.85,2.05) {\textbf{8} ReLU};
% layer 4: output 1 (logistic)
\draw[rounded corners=1pt, fill=white, draw=gray!70, line width=1pt]
  (8.05,1.36) rectangle (9.65,1.74);
\node[font=\scriptsize] at (8.85,1.55) {\textbf{1} logistic};
% inter-layer arrows (main style, >= 1pt)
\draw[main] (8.85,2.86) -- (8.85,2.74);
\draw[main] (8.85,2.36) -- (8.85,2.24);
\draw[main] (8.85,1.86) -- (8.85,1.74);

% ---- decision ----
\node[box, rounded corners=3pt, minimum width=2.5cm, minimum height=1.1cm] at (12.05,3.20) {};
\node[font=\footnotesize\bfseries] at (12.05,3.47) {Risk score};
\node[font=\footnotesize] at (12.05,2.94) {$\hat{y}(r)$};
\node[box, rounded corners=3pt, minimum width=2.5cm, minimum height=1.0cm] at (12.05,2.00) {};
\node[font=\footnotesize\bfseries] at (12.05,2.20) {threshold $\tau$};
\node[box, rounded corners=3pt, minimum width=1.0cm, minimum height=0.7cm,
      fill=failRed!8, draw=failRed!80] at (11.42,0.90) {};
\node[font=\scriptsize\bfseries, text=failRed!90!black] at (11.42,0.90) {FLAG};
\node[box, rounded corners=3pt, minimum width=1.0cm, minimum height=0.7cm,
      fill=passGreen!8, draw=passGreen!80] at (12.79,0.90) {};
\node[font=\scriptsize\bfseries, text=passGreen!90!black] at (12.79,0.90) {accept};
\node[flw] at (12.05,0.36) {5\% FPR target};

% ---- data-flow labels ----
\node[font=\scriptsize\itshape, text=black!85, anchor=west] at (6.48,6.28) {OOF scores};
\node[font=\scriptsize\itshape, text=black!85, anchor=east] at (9.28,4.13) {calibrated $\sigma_i$};

% ---- arrows ----
% input -> sources (fan-out bus)
\draw[main] (2.30,3.25) -- (2.55,3.25);
\draw[bus] (2.55,3.25) -- (2.55,6.05);
\draw[bus] (2.55,3.25) -- (2.55,2.30);
\foreach \y in {6.05,4.80,3.55,2.30} \draw[aux] (2.55,\y) -- (2.90,\y);
% sources -> fusion (converge bus; OOF scores annotated)
\foreach \y in {6.05,4.80,3.55,2.30} \draw[aux] (6.25,\y) -- (6.36,\y);
\draw[bus] (6.36,2.30) -- (6.36,6.05);
\draw[main] (6.36,6.05) -- (8.15,6.05);
% fusion chain
\draw[main] (9.35,5.55) -- (9.35,5.35);
\draw[main] (9.35,4.35) -- (9.35,3.90);
% fusion -> decision
\draw[main] (10.55,3.20) -- (10.80,3.20);
\draw[main] (12.05,2.65) -- (12.05,2.50);
\draw[main] (12.05,1.50) -- (11.42,1.25);
\draw[main] (12.05,1.50) -- (12.79,1.25);
\end{tikzpicture}
\caption{\textbf{The MSRF knowledge-driven fusion framework.} Four mutually
non-overlapping knowledge sources --- harmful-intent semantics (LoRA
probability, d=3), narrative structure (discourse features, d=4), acoustic
prosody (d=5), and scoring uncertainty (judge disagreement, d=1) --- are
extracted from each response, scored out-of-fold under group-aware splits
(OOF scores), calibrated with isotonic regression ($\sigma_i$), and fused by a
225-parameter MLP ($4\to16\to8\to1$) into a discriminative risk score
$\hat{y}(r)$; the operating threshold $\tau$ targets a 5\% FPR operating point
on the deployment frame. The honest boundary behavior of the framework (G2) is
reported in Section~\ref{sec:fusion-boundary}.}
\label{fig:msrf-framework}
\end{figure}

\subsection{Signal families as heterogeneous knowledge sources}
Each signal family is a knowledge source $i \in \{\mathrm{int},
\mathrm{nar}, \mathrm{ac}, \mathrm{unc}\}$ that produces a score
$x_i(r) \in [0,1]$ for a response $r$ (its text or audio rendering of the
same query):
\begin{align}
  x_{\mathrm{int}}(r) &= g_{\mathrm{int}}\!\big(p_{\mathrm{int}}(r),\,
      p_{\mathrm{int}}(r)^2,\, \ln\!\big(1+\ell(r)\big)\big), \label{eq:intent}\\
  x_{\mathrm{nar}}(r) &= g_{\mathrm{nar}}\!\big(\ell(r),\, n_{\mathrm{sent}}(r),\,
      d_{\mathrm{nar}}(r),\, n_{\mathrm{story}}(r)\big), \label{eq:narrative}\\
  x_{\mathrm{ac}}(r) &= g_{\mathrm{ac}}\!\big(f_0(r),\, e(r),\, \delta(r),\,
      \mathrm{zcr}(r),\, r_{\mathrm{sp}}(r)\big), \label{eq:acoustic}\\
  x_{\mathrm{unc}}(r) &= 2\cdot\min\!\big(p_{\mathrm{unc}}(r),\, 1-p_{\mathrm{unc}}(r)\big). \label{eq:uncertainty}
\end{align}
The intent source $x_{\mathrm{int}}$ is a LoRA fine-tuned harmful-intent
classifier over the E2B base whose input features are the intent probability,
its square, and normalized response length (\eqref{eq:intent}); the narrative
source $x_{\mathrm{nar}}$ uses interpretable discourse-structure features ---
length, sentence count, narrative density, and story-word presence
(\eqref{eq:narrative}); the acoustic source $x_{\mathrm{ac}}$ uses prosodic
features (f0, energy, duration, zero-crossing rate, speech/pause ratios) with
mean-imputation-plus-mask when audio is missing (\eqref{eq:acoustic}); and
the uncertainty source $x_{\mathrm{unc}}$ encodes scoring disagreement across
the multi-scorer protocol as $2\cdot\min(p,\,1-p)$, maximal when the scorer
consensus is maximally uncertain (\eqref{eq:uncertainty}). The families are
designed to be non-overlapping: each covers a disjoint aspect of the response,
so the fusion cannot ``double count'' the same evidence through two branches.
We disclose two boundary conditions of the sources: the intent branch was
trained in-distribution (its training rows overlap the evaluation set by
design of the stacking protocol), and the acoustic branch has no measurable
signal at our sample window ($N=32$ attacks). We phrase the latter as
``not measured'', not ``ineffective''.

\subsection{Fusion rule and parameter budget}
Branch scores are calibrated with isotonic regression $\sigma_i$ fit on
training splits to render them comparable across families, then fused by a
two-hidden-layer MLP $\mathcal{F}$:
\begin{align}
  \hat{y}(r) &=
    \mathcal{F}\!\Big(\sigma_1\big(x_{\mathrm{int}}(r)\big),\,
    \sigma_2\big(x_{\mathrm{nar}}(r)\big),\,
    \sigma_3\big(x_{\mathrm{ac}}(r)\big),\,
    \sigma_4\big(x_{\mathrm{unc}}(r)\big)\Big), \label{eq:fusion}\\
  |\mathcal{F}| &= (4\cdot 16 + 16) + (16\cdot 8 + 8) + (8\cdot 1 + 1) = 225. \label{eq:params}
\end{align}
The fused layer is a $4 \rightarrow 16 \rightarrow 8 \rightarrow 1$ ReLU MLP
with 225 trainable parameters (\eqref{eq:params}); the operating threshold
$\tau$ targets 5\% FPR on the deployment frame. The parameter budget is a
deliberate design choice: with 225 parameters the entire decision function is
inspecting three numbers per source, so every contribution to the score can be
traced back to an interpretable branch. This is an auditability property that
multi-billion-parameter guards (ShieldGemma-9B, WildGuard, GradSafe) do not
offer, and one we exploit to localize the honest boundary of the framework in
Section~\ref{sec:fusion-boundary}. This is the knowledge-fusion position of
\citet{shafer1976} and recent evidential fusion work
\citep{guo2024survey,shao2024duallevel,xu2025promethee}: evidence is made
explicit and combined under uncertainty, rather than buried in a black box.

\begin{table}[t]
\centering
\small
\caption{\textbf{The 225-parameter fusion budget.} Layer-by-layer composition
of the fusion MLP $\mathcal{F}$ ($4\to16\to8\to1$). The four inputs are the
isotonic-calibrated branch scores $\sigma_1,\dots,\sigma_4$ (intent, narrative,
acoustic, uncertainty); hidden activations are ReLU and the output is a
logistic layer, matching the scikit-learn \texttt{MLPClassifier} used in the
protocol (\texttt{hidden\_layer\_sizes=(16,8)}, \texttt{activation="relu"}).}
\label{tab:params}
\begin{tabular}{@{}llll@{}}
\toprule
Layer & Shape & Composition & Parameters \\
\midrule
1 & $4\to16$ (ReLU) & $4\times16+16$ & 80 \\
2 & $16\to8$ (ReLU) & $16\times8+8$ & 136 \\
3 & $8\to1$ (logistic) & $8\times1+1$ & 9 \\
\midrule
\multicolumn{3}{r}{\emph{Total trainable}} & \textbf{225} \\
\bottomrule
\end{tabular}
\end{table}

\subsection{Evaluation protocol}
Branch scores are produced \emph{out-of-fold} so that no row contributes its
own label to the score the fusion layer sees. The split is a group-aware
GroupShuffleSplit whose group key is template family $\times$ attack family
(21 groups), so that an entire template$\times$attack group is held out
together and leakage across conditions is prevented. An outer 25\% test frame
never participates in hyperparameter selection, and an inner 20\% validation
split selects the MLP learning rate. The protocol is repeated over five
independent seeds, yielding five independent test frames with the deployment
frame reported as primary. The silver ground-truth label is an
agreement-based subset of the multi-scorer protocol (66.2\% coverage; mean
pairwise agreement 0.8122), used only to train and score the fused detector.

Algorithm~\ref{alg:msrf} summarizes training and inference.

\begin{algorithm}[t]
\caption{MSRF training and inference}
\label{alg:msrf}
\begin{algorithmic}[1]
\Require responses $R$ (text/audio renderings); silver labels $L\subset R$;
        groups $G$ (template family $\times$ attack family)
\Ensure calibrated fusion model $\mathcal{F}$; threshold $\tau$
\Procedure{Train}{$R, L, G$}
  \For{each family $i \in \{\mathrm{int},\mathrm{nar},\mathrm{ac},\mathrm{unc}\}$}
    \State extract features and compute branch score $x_i(r)$ for all $r\in R$
    \State recompute branch scores out-of-fold over group-aware splits of $G$
  \EndFor
  \State fit isotonic calibrators $\sigma_i$ on the training splits
  \State select MLP learning rate on the inner 20\% validation split
  \State fit $\mathcal{F}: \mathbb{R}^4\to[0,1]$ (4$\to$16$\to$8$\to$1) on calibrated branch scores
  \State set $\tau$ to target 5\% FPR on the deployment-frame validation set
  \State \Return $\mathcal{F},\,\{\sigma_i\},\,\tau$
\EndProcedure
\Procedure{Infer}{$r$}
  \State compute $\{x_i(r)\}$; \Return $\mathcal{F}\big(\sigma_1(x_1),\dots,\sigma_4(x_4)\big)$
\EndProcedure
\end{algorithmic}
\end{algorithm}
